% BHU PYQ -- Transcribed & Corrected
% Paper: PHYSE31: Computational Physics / Python Programming (Sessional Exam)
% Exam: B.Sc. Semester III Sessional Examination, 2025 (NEP)
% Category: Skill Enhancement Course (SEC 3)
% Department: Physics

\documentclass[11pt,a4paper]{article}
\usepackage[a4paper,top=2.5cm,bottom=2.5cm,left=2.8cm,right=2.8cm,headheight=14pt]{geometry}
\usepackage{amsmath,amssymb}
\usepackage[T1]{fontenc}
\usepackage[utf8]{inputenc}
\usepackage{lmodern,microtype,fancyhdr,enumitem,hyperref,lastpage,parskip,listings}

\hypersetup{
    colorlinks=true,
    urlcolor=black,
    linkcolor=black,
    pdftitle={PHYSE31: Python Programming (Sessional Exam) -- BHU Sem III 2025 (NEP SEC)}
}

\pagestyle{fancy}
\fancyhf{}
\renewcommand{\headrulewidth}{0.4pt}
\renewcommand{\footrulewidth}{0.4pt}
\fancyhead[L]{\small \textit{Banaras Hindu University}}
\fancyhead[R]{\small \textit{SEC 3: Python Programming (Sessional 2025)}}
\fancyfoot[C]{\small Page \thepage\ of \pageref{LastPage}}

\newcommand{\pts}[1]{\hfill \textit{[#1]}}

\lstset{
  language=Python,
  basicstyle=\ttfamily\small,
  keywordstyle=\bfseries,
  showstringspaces=false,
  frame=single,
  frameround=tttt
}

\begin{document}

\begin{center}
  {\large\bfseries BANARAS HINDU UNIVERSITY}\\[2pt]
  {\normalsize B.Sc. (Hons.) --- Semester III Sessional Examination, 2025 (NEP)}\\[6pt]
  \rule{0.8\linewidth}{0.5pt}\\[6pt]
  {\large\bfseries PHYSICS (SEC 3)}\\[4pt]
  {\large\bfseries Paper Code: PHYSE31 : Computational Physics --- Python Programming (Sessional)}\\[4pt]
  {\normalsize Time: 55 Minutes\hfill Maximum Marks: 20}
\end{center}

\rule{\linewidth}{0.4pt}

\smallskip

\noindent \textit{Instructions: Use the Python programming language to address the following problems.}

\bigskip

\noindent \textbf{Question 1.} Print the following outputs: \pts{1+1+1}
\begin{enumerate}[label=\alph*)]
  \item your name
  \item your roll number
  \item A finite sequence $X$ with elements $a_n = n^2$ where $n = [1, 10]$
\end{enumerate}

\bigskip

\noindent \textbf{Question 2.} Find the mean of the finite sequence $X = \{n^2\}$ where $n = [1, 10]$. \pts{3}\\
\textit{Hint:} $(a_n)_{\text{mean}} = \mu = \frac{a_1 + a_2 + a_3 + \dots + a_N}{N}$

\bigskip

\noindent \textbf{Question 3.} Find the standard deviations of the finite sequence $X = \{n^2\}$ where $n = [1, 10]$. \pts{4}\\
\textit{Hint:} $(a_n)_{\text{std}} = \sigma \equiv \sqrt{\frac{(a_1 - \mu)^2 + (a_2 - \mu)^2 + (a_3 - \mu)^2 + \dots + (a_N - \mu)^2}{N}}$

\bigskip

\noindent \textbf{Question 4.} What is the use of negative indexes? Demonstrate it using an example of your choice? \pts{2}

\bigskip

\noindent \textbf{Question 5.} Write a python program to find the factorial of 6. \pts{3}

\bigskip

\noindent \textbf{Question 6.} Mention True or False for the following statement: \pts{5}
\begin{enumerate}[label=\alph*)]
  \item Python is written in c language.
  \item A hello function can be defined as \texttt{"def Hello():"}.
  \item A library can be imported using \texttt{"import numpy as np1"}.
  \item A numpy array can be created using \texttt{"nums = [1,2,3,4,5]"}.
  \item The output of the \texttt{"print(4 + '3')"} is 7.
\end{enumerate}

\bigskip
\begin{center}
  \textbf{***}
\end{center}

\end{document}
